\subsection{Capacity under demand-proportional allocation (RQ2, H2)}
\label{sec:res-alloc}

\begin{table}[t]
\caption{Equal versus demand-proportional CPU allocation in the local environment: measured peak throughput
(requests/s), the same relative to the monolith measured in the same campaign, concurrency levels at
which a configuration measured above the monolith, and the busiest container at 64 clients with its CPU
utilisation.}
\label{tab:alloc}
\centering\scriptsize
\setlength{\tabcolsep}{3pt}
\begin{tabular}{@{}llrrrrlll@{}}
\toprule
& & \multicolumn{2}{c}{Peak} & \multicolumn{2}{c}{vs.\ A} & & \multicolumn{2}{c}{Busiest container at $N=64$} \\
\cmidrule(lr){3-4}\cmidrule(lr){5-6}\cmidrule(l){8-9}
Workload & Configuration & equal & prop. & equal & prop. & Wins over A & equal & prop. \\
\midrule
\xoinput{generated/tab-alloc.tex}
\bottomrule
\end{tabular}
\end{table}

Table~\ref{tab:alloc} compares the equal split with the demand-proportional allocation of
Section~\ref{sec:treatments}. Relative to the monolith measured alongside, moving the budget to where the demand
was raised the capacity of the decomposed configurations by a factor of \xo{alloc/gainmin} to
\xo{alloc/gainmax}. B reached \xo{prop/peakpct/browse/B}\% of the monolith's peak on browse and
\xo{prop/peakpct/checkout/B}\% on checkout, against \xo{peakpct/browse/B}\% and \xo{peakpct/checkout/B}\% under
the equal split. C reached \xo{prop/peakpct/browse/C}\% and \xo{prop/peakpct/checkout/C}\%, against
\xo{peakpct/browse/C}\% and \xo{peakpct/checkout/C}\%. The monolith's allocation is unchanged, and its peak of
\xo{prop/peak/browse/A} browse requests per second was close to its \xo{peak/browse/A} in the first
campaign. No decomposed configuration matched the monolith at any measured concurrency level, so \textbf{H2 is
falsified under this treatment as well}.

The busiest containers explain why. The monolith was still bound by its database. B and C were still bound by
their edge services, now at the limit of a two- to fourfold larger quota (\texttt{\xo{alloc/edge/browse/C}} in~C:
\xo{alloc/cores/browse/C} cores instead of \xo{alloc/equal/C}). The exception is checkout in~B, where
\texttt{order} was no longer the limit and \texttt{inventory} took its place. Reallocation removed the idle shares
but not the extra work. At 600 browse requests per second, C's application tier needed
\xo{prop/cpums/browse/C/600/app}\,ms of CPU per request against the monolith's \xo{prop/cpums/browse/A/600/app}\,ms.
At that cost even a perfect split of the four application cores serves fewer than half the requests per second
at which the monolith's database saturated.

H1 and H3 remained supported. The p99 ratios to the monolith at the lowest rate were
\xo{prop/h1/min}--\xo{prop/h1/max}, against
\xo{h1/min}--\xo{h1/max} under the equal split. The allocation did not cause this
difference. At these rates no edge container was throttled in either campaign, and the decomposed configurations'
p99 barely changed (browse in~B: \xo{p99/browse/B/50} and \xo{prop/p99/browse/B/50}\,ms). The monolith, whose
allocation is identical in both campaigns, measured slower in the second one: \xo{p99/browse/A/50} against
\xo{prop/p99/browse/A/50}\,ms on browse and \xo{p99/checkout/A/25} against \xo{prop/p99/checkout/A/25}\,ms on
checkout. Low-load tail latency therefore drifts between campaigns on this host by more than the intervals within
one campaign suggest. Section~\ref{sec:threats} returns to this point.
